\section{Introduction}\label{sec:introduction}

How does spatial price competition change when a destination's demand can move transport service toward that destination and away from its rival? Standard spatial-competition models usually treat access as fixed. In many managed transport systems, however, service intensity responds to passenger demand while the total service resource is scarce. A destination that attracts more trips can then receive more service only by drawing that resource away from another direction or zone. Retail pricing affects not only who shops where, but also the access conditions under which the two destinations compete.

We study this feedback in a Hotelling--Bertrand model with two endpoint retail destinations and a third-party transport operator. The operator reallocates a fixed total service resource across the two directions in response to total passenger demand. One direction also carries exogenous background demand, such as commuting or education trips. A retailer's price cut raises its shopping demand, shifts service toward its own direction, withdraws service from the rival direction, and thereby feeds back into shopping demand. Directional background demand makes this shared-resource feedback asymmetric.

The main result is that the feedback can reverse the strategic relation for only one retailer. There exists a nonempty open set of primitives for which the model has a global pure-strategy price Nash equilibrium satisfying
\begin{equation*}
\frac{dBR_L}{dp_R}<0<\frac{dBR_R}{dp_L}.
\end{equation*}
Thus one retailer lowers its price when the rival raises price, while the other retains the conventional upward-sloping response. Strategic asymmetry itself is known \citep{tombak2006}, and mixed strategic relations also arise in transport-network pricing \citep{mchardy2024}. Our contribution is narrower: the sign pattern is generated by downstream demand moving a third party's shared transport resource, and it survives the global retail-price game rather than appearing only in local first-order conditions.

The minimum service obligation has a separate global role. Without a lower service bound, a sufficiently large price-induced demand shift can drive service in the losing direction toward zero and create multiple fulfilled-expectations continuations, a form of demand--frequency feedback familiar in transit economics \citep{baryosef2013}. We show that a minimum service share can be strictly slack at the equilibrium and yet support that equilibrium by truncating only sufficiently large off-path reallocations. For the exact existence witness, the same on-path equilibrium is supported over an intermediate range of nonbinding floor levels. We do not interpret that range as an optimal regulatory rule.

The transport microfoundation also disciplines welfare accounting. Conditional on retail shares, the operator minimizes aggregate waiting cost. The resulting envelope identity makes the directional waiting-cost difference faced by the marginal shopper equal to the derivative of minimized aggregate waiting cost. Retail prices cancel as transfers in welfare, so the decentralized allocation can be compared directly with the constrained-efficient allocation under the same service floor. At the exact witness, the decentralized equilibrium assigns $57.5\%$ of shoppers to the high-background-demand side, versus approximately $70.24\%$ in the same-floor constrained-efficient allocation, reducing real spatial and waiting costs by about $1.22\%$ under the normalization. This is a welfare implication of the transport mechanism, not a claim that envelope arguments or private--social wedges are new in general.

Three nested benchmarks identify the interaction. The opposite-signed local price reactions disappear if service is fixed, if directional background demand is removed, or if service does not respond to shopping demand. The result is also locally robust when waiting disutility is generalized from $w/f$ to $w f^{-\rho}$ around $\rho=1$. The exact parameter vector is a constructive existence device; the economic claim is the open-set mechanism, not the particular witness.

The model is intended for settings in which downstream destinations are large enough to affect discretionary trip demand and service allocation. Directional shuttle or fixed-route operations using deadheading or short-turning provide one interpretation \citep{furth1985,cortes2011}; demand-responsive or shared-mobility fleets repositioned across zones provide another \citep{liuouyang2021}. The operator is a service-planning layer with a fixed effective resource, not a private platform whose fare, fleet size, or network design is modeled. We therefore do not claim that an ordinary small retailer changes public-transit frequency or that the model describes a generic balanced two-way timetable.

The paper sits at the intersection of spatial competition, network externalities, vertical transport markets, capacity allocation, and transit service design. Market-share feedback in Hotelling competition is well established \citep{grilo2001,grivavettas2011,tolottiyepez2020,toshimitsu2026}; upstream transport facilities and service quality have long been studied jointly with downstream competition \citep{bassozhang2007,depalma2018,alvarez2020}; and scarce upstream capacity allocation is known to reshape downstream rivalry \citep{chenli2013}. The remaining claim is the complete transport-specific chain in which downstream demand reallocates one common service resource across rival directions, altering both own and rival access and generating the global price equilibrium above. Section~\ref{sec:literature} develops this positioning.

The remainder of the paper is organized as follows. Section~\ref{sec:model} presents the model and operator allocation. Section~\ref{sec:equilibrium} derives the price-game mechanism and nested benchmarks. Section~\ref{sec:global} establishes the global equilibrium and the off-path role of the service obligation. Section~\ref{sec:welfare} develops the welfare mapping. Section~\ref{sec:robustness} discusses robustness, institutional interpretation, and empirical implications. Section~\ref{sec:literature} positions the paper against the closest literatures, and Section~\ref{sec:conclusion} concludes.
